% Define a subsection
\subsection{Dorado Basecaller}
% Use an acronym

Dorado je pokročilý basecaller vyvinutý spoločnosťou Oxford Nanopore Technologies (ONT) na prekladanie surového elektrického signálu získaného z nanopórového sekvenovania na sekvenciu nukleotidov (A, T, C, G pre DNA alebo A, U, C, G pre RNA).
Je navrhnutý tak, aby poskytoval vysokú presnosť a rýchlosť pri spracovaní sekvenačných dát, čo je kľúčové pre rôzne aplikácie v oblasti genomiky a bioinformatiky.

\subsubsection{Ako spustiť Dorado}

% dorado basecaller hac ../../../../media/marek/EAGET/Bioinformatic/MMC234__202311/ --recursive --max-reads 100 > testCalls.bam
% dorado basecaller fast ../../../../media/marek/EAGET/Bioinformatic/MMC234__202311/ --recursive > fullCall.bam

Pred spustením bolo potrebné stiahnuť model pomocou príkazu:
\begin{verbatim}
dorado download --list
dorado download --model dna_r10.4.1_e8.2_400bps_hac@v4.2.0
\end{verbatim}

% dorado download --list vypíše zoznam dostupných modelov
% dorado download --model <nazov> stiahne konkrétny model

Dorado podporuje viacero módov spracovania:
\begin{itemize}
    \item \texttt{fast} — rýchlejší, ale menej presný
    \item \texttt{hac} — pomalší, ale kvalitnejší (High Accuracy)
    \item \texttt{sup} — najpomalší, najvyššia presnosť (Super Accuracy)
\end{itemize}

Na spustenie basecallera som použil nasledujúci príkaz:
\begin{verbatim}
dorado basecaller hac <adresar_s_pod5> --recursive --max-reads 100 > testCalls.bam
\end{verbatim}

% --recursive je potrebné pridať ak chceme spracovať všetky súbory v adresári a jeho podadresároch
% --max-reads obmedzí počet spracovaných čítaní — vhodné na testovanie na menšom množstve dát

Parameter \texttt{--recursive} je potrebný na spracovanie všetkých súborov v adresári a jeho podadresároch.
Dorado je dôležité spúšťať s prístupom k GPU — bez neho je spracovanie veľmi pomalé. Aj s GPU je spracovanie pomerne pomalé, preto je vhodné pri testovaní používať parameter \texttt{--max-reads}, čo ušetrí veľa času pri práci s miliónmi sekvencií.

% dokumentácia: https://software-docs.nanoporetech.com/dorado/latest/basecaller/basecall_overview/
\cite{DoradoDocs}
